\documentclass[12pt]{article}
\usepackage[utf8]{inputenc}
\usepackage[T1]{fontenc}
\usepackage{hyperref}
\title{Tanglang Language Resources for K-12 Curriculum Development}
\author{Compiled by Claude Sonnet 5.5 (claude-sonnet-5-5) for Andr\'as Kornai}
\date{}
\begin{document}
\maketitle

\section{Introduction}
Tanglang (ISO 639-3 ytl; Glottocode tang1372), also called Tholo, is a Loloish language of Yunnan. It is spoken in eight villages of southern Tai'an Township, Lijiang County. Group population figures vary: one source gives 947, and Wikipedia's infobox lists 2,000 speakers (2009). The language has long been in contact with Bai and Naxi. UNESCO's atlas classes it as Vulnerable. The group is officially counted as Yi, but Bradley suggests it may have been misclassified. Little material was found, and nothing suggests it is taught in schools.

\section{Primary Educational Resources}
No primers, textbooks or school programmes were found.

\section{Linguistic and Reference Materials}
\begin{itemize}
\item Wikipedia, Tholo (Tanglang) language: \url{https://en.wikipedia.org/wiki/Tholo_language}
\item Endangered Languages Project entry: \url{https://endangeredlanguages.com/lang/5175}
\item Verbix Endangered Languages, Tanglang: \url{https://docs.verbix.com/EndangeredLanguages/Tanglang}
\item Glottolog reference record, David Bradley 2004: \url{https://glottolog.org/resource/reference/id/547662}
\end{itemize}

\section{Religious and Community Resources}
\begin{itemize}
\item Joshua Project, Tanglang in China: \url{https://joshuaproject.net/people_groups/18695/CH}
\item Asia Harvest people-group profile: \url{https://people-groups.asiaharvest.org/China/chinaPeoples/T/Tanglang.pdf}
\end{itemize}
No community-run language organization was found.

\section{Digital and Audio Resources}
No audio, video, corpus or software resources for Tanglang were found.

\section{Curriculum Development Context}
The language is spoken by a small group in eight villages, with no known writing system or teaching materials. A curriculum would have to start with documentation and orthography work with the community. The language is not known to be taught in schools.
\end{document}